% part: intuitionistic-logic
% chapter: introduction
% section: bhk-interpretation

\documentclass[../../../include/open-logic-section]{subfiles}

\begin{document}

\olfileid{int}{int}{bhk}

\olsection{براوَر--ہائٹنگ--کولموگوروف دی تعبیر}

\begin{editorial}
  بی ایچ کے تعبیر نال وجدانیت پسند قضیاں دے اعتبار دے ثبوت سمجھنے
  اوکھے نیں۔ اوہناں نوں ہور صاف سمجھاؤن دی لوڑ اے۔
\end{editorial}

وجدانیت پسند رابطیاں دی اک غیر رسمی تعمیری تعبیر اے، جیہنوں عام
طور اُتے براوَر، ہائٹنگ تے کولموگوروف دی تعبیر آکھدے نیں۔ ایہ
«تعمیر» دا خیال ورتدی اے، جیہنوں تسی تعمیری ثبوت سمجھ سکدے او۔
(بی ایچ کے تعبیر وچ اسیں «ثبوت» نہیں آکھدے، تاں جو ایس نوں رسمی
!!{derivation} نظام دے !!a{derivation} نال رلا نہ دیئیے۔) بی ایچ
کے تعبیر ایس سِدھے خیال اُتے وجدانیت پسند رابطیاں دے معنی سمجھاؤندی اے۔

\begin{enumerate}
\item اسیں منّ لیندے آں کہ سانوں پتا اے کہ ایٹمی بیان دی تعمیر
  کیہڑی شے اے۔
\item $!A_1 \land !A_2$ دی تعمیر اک جوڑی $\tuple{M_1, M_2}$ اے،
  جتھے $M_1$،~$!A_1$ دی تعمیر اے تے $M_2$،~$!A_2$ دی تعمیر اے۔
\item $!A_1 \lor !A_2$ دی تعمیر اک جوڑی $\tuple{s, M}$ اے،
  جتھے یا $s$،~$1$ اے تے $M$،~$!A_1$ دی تعمیر اے، یا $s$،~$2$
  اے تے $M$،~$!A_2$ دی تعمیر اے۔
\item $!A \lif !B$ دی تعمیر اوہ فنکشن اے جیہڑا~$!A$ دی تعمیر
  نوں~$!B$ دی تعمیر وچ بدل دیندا اے۔
\item ~$\lfalse$ (محال) دی کوئی تعمیر نہیں۔
\item $\lnot !A$ دی تعریف $!A \lif \lfalse$ دے ہم معنی وجوں ہوندی اے۔
  یعنی $\lnot !A$ دی تعمیر اوہ فنکشن اے جیہڑا~$!A$ دی تعمیر
  نوں~$\lfalse$ دی تعمیر وچ بدل دیندا اے۔
\end{enumerate}

\begin{ex}
مثال وجوں $\lnot \lfalse$ لو۔ ایہدی تعمیر ایسا فنکشن اے جیہڑا
$\lfalse$ دی کوئی وی تعمیر لے کے جواب وچ~$\lfalse$ دی تعمیر دیندا
اے۔ صاف اے کہ شناختی فنکشن~$\Id{}$ ایہو جیہی تعمیر اے۔ جے~$M$،
~$\lfalse$ دی تعمیر اے، تاں $\Id{}(M) = M$ توں~$\lfalse$ دی تعمیر ملدی اے۔
\end{ex}

عام طور اُتے $\lnot !A$ دا مطلب اے کہ $!A$ دی تعمیر ناممکن اے۔

\begin{ex}
$!A \lif \lnot \lnot !A$ ثابت کریے، کسے وی قضیے~$!A$ لئی، یعنی
$!A \lif ((!A \lif \lfalse) \lif \lfalse)$۔ تعمیر ایسا فنکشن~$f$
ہونا چاہیدا اے جیہڑا تعمیر~$M$، $!A$ دی، لے کے تعمیر~$f(M)$
دیندا اے، $(!A \lif \lfalse) \lif \lfalse$ دی۔~$f$،
$(!A \lif \lfalse) \lif \lfalse$ دی تعمیر ایویں بناؤندا اے۔
سانوں ایسا فنکشن $g$ متعین کرنا پئے گا جیہڑا تعمیر~$h$،
$!A \lif \lfalse$ دی، لے کے~$\lfalse$ دی تعمیر دیوے۔ $g$ نوں
ایویں متعین کر سکدے آں۔ ملی ہوئی~$h$ نوں تعمیر~$M$ اُتے لاؤ، جیہڑی
$!A$ دی اے (تے سانوں پہلاں ملی سی)۔ $h(M)$، $h$ توں ملدا جواب،
$\lfalse$ دی تعمیر اے۔ ایس لئی $f(M)(h) = h(M)$،~$\lfalse$ دی تعمیر اے
جے $M$،~$!A$ دی تعمیر اے۔
\end{ex}

\begin{ex}
$\lnot(!A \land \lnot !A)$ دی تعمیر دیئیے، یعنی $(!A
\land (!A \lif \lfalse)) \lif \lfalse$ دی۔ ایہ ایسا فنکشن~$f$
اے جیہڑا تعمیر~$M$، $!A \land (!A \lif \lfalse)$ دی، لے کے
~$\lfalse$ دی تعمیر دیندا اے۔ عطف $!B_1 \land !B_2$ دی تعمیر اک
جوڑی $\tuple{N_1, N_2}$ اے، جتھے $N_1$،~$!B_1$ دی تعمیر اے
تے $N_2$،~$!B_2$ دی تعمیر اے۔ اسیں فنکشن $p_1$ تے $p_2$ متعین
کر سکدے آں جیہڑے $!B_1 \land !B_2$ دی تعمیر توں واری واری
$!B_1$ تے $!B_2$ دیاں تعمیراں کڈھدے نیں:
\begin{align*}
  p_1(\tuple{N_1, N_2}) &= N_1\\
  p_2(\tuple{N_1, N_2}) &= N_2
\end{align*}
$f$ ایویں کم کردا اے۔ پہلاں $p_1$ نوں ملی ہوئی~$M$ اُتے لاؤندا اے۔
ایس توں~$!A$ دی تعمیر ملدی اے۔ فیر $p_2$ نوں~$M$ اُتے لاؤندا اے۔
ایس توں~$!A \lif \lfalse$ دی تعمیر ملدی اے۔ ایہ تعمیر آپ اک فنکشن
~$p_2(M)$ اے جیہڑا~$!A$ دی تعمیر لے کے~$\lfalse$ دی تعمیر دیندا اے۔
دوجے لفظاں وچ، جے اسیں $p_2(M)$ نوں $p_1(M)$ اُتے لائیے، تاں
~$\lfalse$ دی تعمیر ملدی اے۔ سو اسیں $f(M) = p_2(M)(p_1(M))$
متعین کر سکدے آں۔
\end{ex}

\begin{ex}
$((!A \land !B) \lif !C) \lif (!A \lif
(!B \lif !C))$ دی تعمیر دیئیے، یعنی ایسا فنکشن $f$ جیہڑا تعمیر~$g$،
$(!A \land !B) \lif !C$ دی، نوں $(!A \lif (!B \lif
!C))$ دی تعمیر وچ بدل دے۔ تعمیر $g$ آپ اک فنکشن اے
(جیہڑا $!A \land !B$ دیاں تعمیراں توں~$C$ دیاں تعمیراں دیندا اے)۔
تے جواب $f(g)$ اک فنکشن~$h_g$ اے، جیہڑا $!A$ دیاں تعمیراں لے کے
ایسے فنکشن دیندا اے جیہڑے~$!B$ دیاں تعمیراں توں~$!C$ دیاں تعمیراں دیندے نیں۔

ایہ گل سمجھنی اوکھی اے۔ سانوں اک فنکشن $h_g$ بناؤنا پئے گا جیہڑا
$f$ نوں~$g$ دین تے جواب وچ ملے۔~$h_g$ نوں دتی جاندی شے تعمیر~$M$
اے،~$!A$ دی۔ $h_g(M)$ توں ملن والا جواب ایسا فنکشن~$k_M$
ہونا چاہیدا اے جیہڑا تعمیراں~$N$،~$!B$ دیاں، توں~$!C$ دیاں
تعمیراں دیوے۔ منّ لو $k_{g,M}(N) = g(\tuple{M, N})$۔ یاد رکھو
کہ $\tuple{M,
  N}$،~$!A \land !B$ دی تعمیر اے۔ سو $k_{g,M}$،
$!B \lif !C$ دی تعمیر اے۔ ایہ تعمیراں~$N$،~$!B$ دیاں، نوں
~$!C$ دیاں تعمیراں وچ بدلدا اے۔ ہُن $h_g(M) = k_{g,M}$ رکھو۔
ایہ فنکشن تعمیراں~$M$،~$!A$ دیاں، نوں تعمیراں $k_{g,M}$،
$!B \lif !C$ دیاں، وچ بدلدا اے۔ ہُن $f(g) = h_g$ رکھو۔
ایہ فنکشن تعمیراں $g$، $(!A \land !B) \lif !C$ دیاں، نوں
$!A \lif (!B \lif !C)$ دیاں تعمیراں وچ بدلدا اے۔ لو، ہو گیا!{}
\end{ex}

بیان $!A \lor \lnot !A$ نوں تیسرے امکان دے اخراج دا اصول آکھدے آں۔
ایہ کجھ خاص $!A$ لئی ثابت کر سکدے آں (مثلاً $\lfalse \lor
\lnot \lfalse$)، پر عام صورت وچ نہیں۔ وجہ ایہ اے کہ وجدانیت پسند
فصل لئی دونہاں وچوں اک جز دی تعمیر چاہیدی اے۔ پر کجھ بیان ہن نہ
ثابت ہو سکدے نیں نہ ردّ (مثلاً گولڈباخ دا قیاس)۔ تیسرے امکان دے
اخراج دے اصول نوں ردّ وی نہیں کر سکدے، یعنی $\lnot \lnot (!A
\lor \lnot !A)$ لاگو اے۔

\begin{ex}
$\lnot \lnot (!A \lor \lnot !A)$ ثابت کرن لئی سانوں ایسا فنکشن~$f$
چاہیدا اے جیہڑا $\lnot(!A \lor \lnot !A)$ دی تعمیر، یعنی $(!A
\lor (!A \lif \lfalse)) \lif \lfalse$ دی تعمیر، نوں~$\lfalse$
دی تعمیر وچ بدلے۔ دوجے لفظاں وچ سانوں ایسا فنکشن $f$ چاہیدا اے
جیہدے لئی $f(g)$،~$\lfalse$ دی تعمیر اے جے $g$،
$\lnot(!A \lor \lnot !A)$ دی تعمیر اے۔

منّ لو $g$، $\lnot(!A \lor \lnot !A)$ دی تعمیر اے، یعنی ایسا فنکشن
جیہڑا $!A \lor \lnot !A$ دی تعمیر نوں~$\lfalse$ دی تعمیر وچ
بدلدا اے۔ $!A \lor \lnot !A$ دی تعمیر اک جوڑی $\tuple{s, M}$ اے
جتھے یا $s = 1$ اے تے $M$، $!A$ دی تعمیر اے، یا $s = 2$ اے تے
$M$، $\lnot !A$ دی تعمیر اے۔ منّ لو $h_1$ اوہ فنکشن اے جیہڑا تعمیر~$M_1$،
$!A$ دی، نوں $!A \lor \lnot !A$ دی تعمیر وچ بدلدا اے۔ ایہ $M_1$ نوں
$\tuple{1, M_2}$ دیندا اے۔ تے منّ لو $h_2$ اوہ فنکشن اے جیہڑا
تعمیر~$M_2$، $\lnot !A$ دی، نوں $!A \lor \lnot !A$ دی تعمیر وچ
بدلدا اے۔ ایہ $M_2$ نوں $\tuple{2, M_2}$ دیندا اے۔

$k$ نوں $\comp{h_1}{g}$ رکھو۔ ایہ ایسا فنکشن اے جیہڑا~$!A$ دی
تعمیر لے کے~$\lfalse$ دی تعمیر دیندا اے، یعنی ایہ~$!A \lif \lfalse$
یا $\lnot !A$ دی تعمیر اے۔ ہُن $l$ نوں $\comp{h_2}{g}$ رکھو۔
ایہ ایسا فنکشن اے جیہڑا $\lnot !A$ دی تعمیر لے کے~$\lfalse$ دی
تعمیر دیندا اے۔ $k$، $\lnot !A$ دی تعمیر اے، ایس لئی $l(k)$،
~$\lfalse$ دی تعمیر اے۔

سو اسیں ایہ دس دتا اے کہ تعمیر~$g$، $\lnot(!A \lor \lnot !A)$
دی، نوں~$\lfalse$ دی تعمیر وچ کِویں بدلیے۔ یعنی فنکشن $f$، جیہڑا
تعمیر~$g$، $\lnot(!A \lor \lnot !A)$ دی، نوں تعمیر $l(k)$،
~$\lfalse$ دی، وچ بدلدا اے، اوہ $\lnot\lnot(!A \lor \lnot !A)$
دی تعمیر اے۔
\end{ex}

تسی ویکھ سکدے او کہ بی ایچ کے تعبیر نال !!{formula}s دا وجدانیت
پسند اعتبار وکھانا چھیتی ای لمبا تے الجھواں ہو جاندا اے۔ چنگی گل اے
کہ وجدانیت پسند منطق لئی ہور چنگے !!{derivation} نظام تے ہور عین
معنیاتی تعبیراں موجود نیں۔
\end{document}
